\chapter{Introduction}
\label{ch:introduction}

As a graduate student paying one-hundred dollars per month across several AI services, I watched peers go without AI tools that had become part of my daily workflow. The technology itself is expensive to run, and the delivery model is not designed for easy access. Syed et al.~\cite{syed2025aiaas} documented this systems problem: the dominant pattern requires users to create an account, supply payment credentials, obtain an API key, and manually configure that key into every application. A student seeking three capabilities from three providers must manage three subscriptions, three keys, and three billing relationships. These barriers create measurable disparities; Bassignana et al.~\cite{bassignana2025aigap} and Gabriel~\cite{gabriel2024educational} showed that users from lower socioeconomic backgrounds interact with language technologies less frequently and derive less benefit. I wanted to provision AI services to everyone on my campus the way universities already provision other shared resources. A student who joins the eduroam network gains access to printers, file shares, and licensed software without creating accounts or entering credentials. There is no zero-configuration protocol advertises, discovers, and connects to AI service endpoints.

I define \emph{Saturn} as the process of advertising AI service API endpoints via mDNS packets on a local network, discovering those endpoints through standard DNS-SD queries, and transferring data between client and service using the discovered connection parameters. A \emph{Saturn service} is any AI endpoint advertised via the \texttt{\_saturn.\_tcp.\allowbreak{}local.} service type. Saturn extends the existing TXT record mechanism to carry endpoint URLs, API types, and authentication credentials without protocol modifications. I created six implementations that consume the same service type with no shared discovery code, no bridging layer, and no user configuration. Python, Rust, and TypeScript clients each resolve Saturn independently. This cross-platform interoperability establishes the first claim (\argref{C1}): zero-configuration network protocols can provision AI services without end-user configuration.

The second claim addresses whether this model actually reduces work (\argref{C2}). A cognitive walkthrough comparing Saturn against traditional per-user API provisioning shows that Saturn reduces combined configuration steps by 53\%, with application developers experiencing a 79\% reduction. The gain is not free: Saturn centralizes all setup in one administrator, who takes on more steps than a single traditional user would. But that one operator replaces the configuration burden of every consumer on the network. 

Zero-configuration access requires a security trade-off. Saturn broadcasts service metadata to every device on the network, and requiring per-device authentication would reintroduce the credential burden the protocol exists to eliminate. I chose to document and bound the exposure instead. The mDNS threat models that Kaiser and Waldvogel~\cite{kaiser2014privacy} established map the attack surface, and ephemeral credentials with ten-minute expiry windows limit what a compromised key can do. I analyze these trade-offs against the static API key model Saturn displaces in the security section of Chapter~\ref{ch:discussion}.

Together, these contributions argue that AI access can and should be provisioned at the network level. Saturn demonstrates the technical feasibility of this model, quantifies its usability benefits, and maps the security surface that any future deployment must address.
